\documentclass[11pt,a4paper]{letter}
\usepackage[a4paper,margin=1.8cm]{geometry}
\usepackage[T1]{fontenc}
\usepackage[utf8]{inputenc}
\usepackage{siunitx}
\usepackage{amsmath}
\usepackage[colorlinks=true,linkcolor=blue,urlcolor=blue]{hyperref}
\sisetup{detect-all}
\setlength{\parskip}{0.4em}
\setlength{\parindent}{0pt}

\address{%
Serge Guy Nana Engo\\
Department of Physics, Faculty of Science\\
University of Yaound\'e I\\
P.O. Box 812, Yaound\'e, Cameroon\\
\texttt{serge.nana-engo@facsciences-uy1.cm}}
\signature{Serge Guy Nana Engo\\on behalf of all authors}
\date{\today}

\begin{document}
\begin{letter}{The Editors\\\textit{Chaos, Solitons \& Fractals}}
\opening{Dear Editors,}

We submit the manuscript entitled ``Synthetic-flux control of hyperchaos order in a dissipative optomechanical model'' for consideration as a full-length research article in \textit{Chaos, Solitons \& Fractals}.

The work addresses a focused question in dissipative nonlinear dynamics: whether a gauge-invariant synthetic-flux phase can act as a reproducible control coordinate for the \emph{order} of a drive- and coupling-gated hyperchaos transition. Within a normalized six-dimensional optomechanical model (one cavity + two mechanical resonators), the flux phase selects up to four simultaneously positive Lyapunov directions---beyond the order of two reported for quadratic single-mode optomechanics (Halef and Shomroni, \textit{Phys.\ Rev.\ Lett.}\ \textbf{135}, 257201 (2025)). The integer count $n_+=4$ is phase- and threshold-sensitive at the extrema. The route proceeds via a supercritical Neimark--Sacker bifurcation at $E^{*}=\num{1.060}$ ($\theta=0$) and $E^{*}=\num{0.975}$ ($\theta=\pi/2$), with re-stabilization at $E=\num{2.0}$ and finite-time counts of up to three positive exponents at $E=\num{4.0}$ and up to four at $E=\num{8.0}$.

Confidence in the classification rests on a phase-consistent Floquet--Lyapunov protocol (monodromy multipliers, dissipative volume balance $\sum\lambda_i=\operatorname{Tr}(J)=-1.04$, \num{432}-run classification ensemble spanning twelve seeds, twelve phases, and three drives), orbit continuation and window-convergence tests. As a secondary geometric clarification, the matched-resource force-sensing protocol yields only a modest enhancement on the stable orbit ($\num{1.039}\pm\num{0.014}$) and a null result on the chaotic attractor ($\num{0.54}\times$ at the resolved $E=\num{4}$ points); this null is specific to the matched-resource protocol and is presented as a geometric bound, not as a general verdict on chaotic sensing.

All results are strictly model-level: the explored sector lies $\num{2542}\times$ beyond anchored silicon couplings, with robust hyperchaos onset near $g_1/\omega_m\approx\num{0.18}$. Closing that gap requires ultrasonic characterization of $\omega_2$, a measured inter-resonator hopping $J_m$, and a genuine gigahertz platform. Four limitations are stated: (i)~the classification ensemble is not a uniqueness proof and $n_+=4$ remains threshold- and seed-sensitive at marginal phases; (ii)~the null sensing result is specific to the matched-resource constraint; (iii)~Fock/Wigner checks do not bound quantum corrections to the Lyapunov spectrum or Fisher information; (iv)~the strong-coupling sector lies far above present experimental anchors. We make no claim of quantum advantage or experimental validation.

The manuscript fits the journal's scope on high-dimensional nonlinear dynamics, bifurcation routes and their physical applications. The control studied is open-loop and parametric. Every numerical claim is produced by deterministic computation. This submission is made under the journal's subscription route.

This work is original, has not been published previously, and is not under consideration elsewhere. All authors have approved the manuscript and declare no competing interests. No funding was received. Thank you for your consideration.

\closing{Sincerely,}
\end{letter}
\end{document}
